% GENERATED FILE. Edit canonical lessons, then run npm run generate:books.
% canonical-source-hash: fnv1a64:db425136e02195fa
% canonical-lessons: PT-C267-se, PT-C267-se-chove, PT-C267-depende, PT-C267-revisao-condicoes

\chapter{Se, and Depende --- If, and It Depends}
\label{ch:pt-if-and-it-depends}
\hlchaptermodality{267}

\begin{chapteropening}
\textbf{By the end of this chapter:} \emph{I can make a real condition with se and the present, and answer with depende when the condition is still open.}

Complete PT-C267-revisao-condicoes: plans with a branch in them, from decided to open.
\end{chapteropening}

\section[se]{se — if}

\label{lesson:PT-C267-se}



\begin{quote}
Say \textquotedblleft{}I don't know\textquotedblright{}. (\emph{Não sei.}) Then \textquotedblleft{}I have been studying Portuguese for three years\textquotedblright{}. (\emph{Há três anos que estudo português.})

Every sentence so far has said what is, or what was. This chapter says what follows \textbf{if} something else is true, and it begins with two letters.
\end{quote}

\begin{grammarlens}[title={Grammar Lens: se opens a condition}]
\textbf{se} — \textquotedblleft{}if\textquotedblright{}.

\begin{quote}
\textbf{Se} tens fome, come. — If you're hungry, eat. \textbf{Se} tens frio, fecha a janela. — If you're cold, shut the window.
\end{quote}

Whatever follows \emph{se} is the \textbf{condition}; the rest of the sentence is what follows from it. The condition can come first, with a comma, or second, without one: \emph{Fecha a janela \textbf{se} tens frio.} Notice that Portuguese \textbf{has} hunger and cold, \emph{tens fome}, \emph{tens frio}, where English \textbf{is} hungry and cold.

\textbf{se} also means \textquotedblleft{}whether\textquotedblright{}, after a verb of knowing or asking, and here your \emph{não sei} comes in useful:

\begin{quote}
\textbf{Não sei se} é verdade. — I don't know if it's true.
\end{quote}
\end{grammarlens}

\begin{cousinweb}
You have met \textbf{se} before, in \emph{como se chama} and \emph{levantar-se}. That one means \textquotedblleft{}oneself\textquotedblright{}, and it is a different word that happens to be spelled the same way. EtymDB 2.1 keeps the two apart: \textbf{se}, \textquotedblleft{}if\textquotedblright{}, is inherited from Latin \emph{sī}, \textquotedblleft{}if\textquotedblright{} (Lewis \& Short, under \emph{si}), whose vowel opened to \emph{e}; and \textbf{se}, \textquotedblleft{}oneself\textquotedblright{}, is inherited from Latin \emph{sē}, \textquotedblleft{}himself, herself, itself\textquotedblright{}.

The ear tells them apart by their place. The \textquotedblleft{}oneself\textquotedblright{} \emph{se} clings to a verb: \emph{chama-se}, \emph{levanta-se}, \emph{como se chama}. The \textquotedblleft{}if\textquotedblright{} \emph{se} opens a whole clause.
\end{cousinweb}

\subsection*{Your turn}
Say these aloud:

\begin{itemize}
\raggedright
  \item \emph{Se tens fome, come.}
  \item the same, condition second — \emph{Come, se tens fome.}
  \item \emph{Se tens frio, fecha a janela.}
  \item \emph{Não sei se é verdade.}
\end{itemize}

\begin{tcolorbox}[breakable,colback=teal!4,colframe=teal!35!black,title={Before you move on}]
Which half of \emph{se tens fome, come} is the condition? (\emph{\textbf{Se tens fome.}}) What does \emph{se} mean after \emph{não sei}? (\textbf{Whether}, \textquotedblleft{}if\textquotedblright{}.) Which Latin words are behind the two kinds of \emph{se}? (\emph{\textbf{Sī}}, \textquotedblleft{}if\textquotedblright{}, and \emph{\textbf{sē}}, \textquotedblleft{}oneself\textquotedblright{}.)
\end{tcolorbox}
\section[se chove, não saio]{se chove, não saio — the condition that really happens}

\label{lesson:PT-C267-se-chove}



\begin{quote}
Say \textquotedblleft{}if\textquotedblright{}. (\emph{Se.}) Say \textquotedblleft{}it was raining\textquotedblright{}. (\emph{Chovia.}) Now the present: \textquotedblleft{}it rains\textquotedblright{}. (\emph{Chove.})
\end{quote}

\begin{grammarlens}[title={Grammar Lens: both halves in the present}]
\begin{quote}
\textbf{Se chove, não saio.} — If it rains, I don't go out. \textbf{Se o comboio é caro, apanho o autocarro.} — If the train is expensive, I take the bus.
\end{quote}

This is a \textbf{real} condition: something true right now, or true whenever it happens. Portuguese puts both halves in the plain present, as English does, so there is nothing new to build.

The result half can also be advice, with the \emph{devias} of chapter two hundred and sixty-two:

\begin{quote}
\emph{Se estás constipado, \textbf{devias} descansar.} — If you've got a cold, you should rest.
\end{quote}
\end{grammarlens}

\begin{culture}
A real condition is a plan with a branch in it. Someone who takes the train to work every day lives by two of them:

\begin{quote}
\emph{Se perco o comboio, apanho o autocarro. Se perco o autocarro, volto para casa.}
\end{quote}

Compare \emph{por isso}. \emph{Chove, por isso não saio} says it \textbf{is} raining, and here is the consequence. \emph{Se chove, não saio} says it \textbf{might}, and here is what happens if so. Same consequence; the second one has not decided yet whether the cause is true.
\end{culture}

\begin{grammarlens}[title={Grammar Lens: one occasion still to come}]
The present is right for a habit, or for what holds now. For \textbf{one} particular occasion still to come, standard Portuguese bends the verb after \emph{se} into a form of its own, the \textbf{future subjunctive}:

\begin{quote}
\emph{Se \textbf{chover} amanhã, não saio.} — If it rains tomorrow, I won't go out.
\end{quote}

\emph{Chover}, not \emph{chove}. For a regular verb it sounds just like the infinitive; irregular verbs have forms of their own, which a later chapter builds. For now, \textbf{recognise} it. Brazil and Portugal share it.
\end{grammarlens}

\begin{cousinweb}
\textbf{chover} is inherited from spoken Latin \emph{plovere} (EtymDB 2.1), a reshaping of classical \emph{pluere}, \textquotedblleft{}to rain\textquotedblright{} (Lewis \& Short, under \emph{pluo}). Latin \emph{pl-} became Portuguese \emph{ch-}: \textbf{chuva} from \emph{pluvia} and \textbf{chegar} from \emph{plicāre} (both EtymDB 2.1), and \textbf{cheio}, \textquotedblleft{}full\textquotedblright{}, usually traced to \emph{plēnus} (Lewis \& Short). Four words, one sound change.
\end{cousinweb}

\subsection*{Your turn}
Say these aloud:

\begin{itemize}
\raggedright
  \item \emph{Se chove, não saio.}
  \item the branch — \emph{Se perco o comboio, apanho o autocarro.}
  \item advice in the result — \emph{Se estás constipado, devias descansar.}
  \item hear the special form — \emph{Se chover amanhã, não saio.}
\end{itemize}

\begin{tcolorbox}[breakable,colback=teal!4,colframe=teal!35!black,title={Before you move on}]
Which tense goes in each half of a real condition? (The \textbf{present}, in both.) What is the difference between \emph{chove, por isso não saio} and \emph{se chove, não saio}? (The first says it \textbf{is} raining; the second has \textbf{not decided}.) Which other words show Latin \emph{pl-} becoming \emph{ch-}? (\emph{\textbf{Chuva}}, \emph{\textbf{cheio}}, \emph{\textbf{chegar}}.)
\end{tcolorbox}
\section[depende]{depende — the condition you do not have to finish}

\label{lesson:PT-C267-depende}



\begin{quote}
Say \textquotedblleft{}if it rains, I don't go out\textquotedblright{}. (\emph{Se chove, não saio.})

That sentence names its condition. Sometimes you want the condition without saying what it is.
\end{quote}

\subsection*{You'll want to know: depende}
\textbf{depende} — \textquotedblleft{}it depends\textquotedblright{}.

\begin{quote}
— \emph{Vais à festa amanhã?} — \emph{Depende.} — Are you going to the party tomorrow? — It depends.
\end{quote}

You have \textbf{depender}, \textquotedblleft{}to depend\textquotedblright{}. \emph{Depende} is its \textquotedblleft{}it\textquotedblright{} form, and on its own it is a complete answer: it says a condition exists and declines to say which. That is exactly why it is so useful for a question you are not ready to answer.

\begin{grammarlens}[title={Grammar Lens: depende de}]
Add \textbf{de} and you can name the condition after all:

\begin{quote}
\textbf{Depende de} ti. — It depends on you. \textbf{Depende do} comboio. — It depends on the train. \textbf{Depende do} dinheiro. — It depends on the money.
\end{quote}

\emph{De} and \emph{o} weld into \textbf{do}, and \emph{de} and \emph{a} into \textbf{da}: \emph{depende da hora}. English says \emph{depend \textbf{on}}; Portuguese says \emph{depende \textbf{de}}, \textquotedblleft{}hangs \textbf{from}\textquotedblright{}, which is the older picture.

Or name it with a whole \emph{se}: \emph{Depende. Se chove, não vou.} The \emph{depende} announces a condition; the \emph{se} sentence states it.
\end{grammarlens}

\begin{cousinweb}
\textbf{depender} comes from Latin \emph{dēpendēre}, \textquotedblleft{}to hang from, to hang down\textquotedblright{} (EtymDB 2.1; Lewis \& Short, under \emph{dependeo}): \emph{dē}, \textquotedblleft{}down from\textquotedblright{}, and \emph{pendēre}, \textquotedblleft{}to hang\textquotedblright{}. To depend on something is to \textbf{hang from} it, and the answer swings until the condition settles. Chapter nineteen hung \textbf{pensar} from the same family, weighing by hanging; English \textbf{pendant} and \textbf{pending} belong to it too.
\end{cousinweb}

\subsection*{Your turn}
Say these aloud:

\begin{itemize}
\raggedright
  \item answer \emph{Vais à festa amanhã?} — \emph{Depende.}
  \item name it — \emph{Depende do comboio.}
  \item name it with a whole condition — \emph{Depende. Se chove, não vou.}
  \item hand it back — \emph{Depende de ti. O que achas?}
\end{itemize}

\begin{tcolorbox}[breakable,colback=teal!4,colframe=teal!35!black,title={Before you move on}]
What does \emph{depende} say, and what does it not say? (That there \textbf{is} a condition — not \textbf{which}.) What do you add to name it? (\emph{\textbf{De}}: \emph{depende de ti}.) What is hanging, in \emph{dēpendēre}? (\textbf{The answer}, from the condition.)
\end{tcolorbox}
\section[(review)]{Review — plans with a branch in them}

\label{lesson:PT-C267-revisao-condicoes}



\begin{quote}
Someone asks \emph{Vais à festa amanhã?} Give the answer that decides nothing. (\emph{Depende.})
\end{quote}

\begin{grammarlens}[title={What you've built}]
Three ways now to talk about a cause, and they differ in how sure you are.

\begin{quote}
\emph{Chove, \textbf{por isso} não saio.} — It is raining: decided. \emph{\textbf{Se} chove, não saio.} — It might: the branch is ready. \emph{\textbf{Depende.}} — Not even the branch is named yet.
\end{quote}

Put them together at a ticket window and you have a whole conversation:

\begin{quote}
— \emph{Queria um bilhete para Lisboa.} — \emph{Ida e volta?} — \emph{Depende. Se volto hoje, ida e volta. Se fico, só de ida.}
\end{quote}

\emph{Fico} is \textquotedblleft{}I stay\textquotedblright{}, from the \emph{ficar} that tells where places are: a place stays where it is.
\end{grammarlens}

\subsection*{Your turn}
Say these aloud:

\begin{itemize}
\raggedright
  \item decided — \emph{Chove, então não saio.}
  \item a branch — \emph{Se chove, não saio.}
  \item not even the branch — \emph{Depende.} — then name what it hangs from — \emph{Depende do comboio.}
  \item at the window — \emph{Depende. Se volto hoje, ida e volta. Se fico, só de ida.}
  \item one occasion still to come — \emph{Se chover amanhã, apanho o autocarro, porque é mais rápido do que ir a pé.} — \emph{a pé}, on foot
\end{itemize}

\begin{tcolorbox}[breakable,colback=teal!4,colframe=teal!35!black,title={Before you move on}]
Order the three from most sure to least. (\emph{\textbf{Por isso}}, then \emph{\textbf{se}}, then \emph{\textbf{depende}}.) Which tense goes in both halves of \emph{se chove, não saio}? (The \textbf{present}.) Which Latin word is the \emph{se} of a condition? (\emph{\textbf{Sī}}, \textquotedblleft{}if\textquotedblright{}.)
\end{tcolorbox}
